% 章の導入
\newenvironment{chapterabstract}{%
  \begin{chapterabstractbox}
  {\sffamily\bfseries\color{MutedBlue}はじめに}\par\smallskip
}{%
  \end{chapterabstractbox}
}

\newcommand{\BookIndex}[2]{\index{#1@#2}}
\newcommand{\term}[1]{\textbf{#1}\index{#1}}
\newcommand{\guideterm}[1]{\textbf{#1}}
\newcommand{\code}[1]{\texttt{#1}}

% 生成画像へ正確な日本語ラベルを図中組版します。
% image-genは構造と余白を作り、文字はTeXで重ねることで誤字と低解像度化を防ぎます。
\newcommand{\BookAnnotatedImage}[3]{%
  \begin{tikzpicture}
    \node[inner sep=0pt,outer sep=0pt] (BookImage) {%
      \includegraphics[width=#1]{assets/#2}%
    };
    \if\relax\detokenize{#3}\relax\else
      \begin{scope}[
        shift={(BookImage.south west)},
        x={($(BookImage.south east)-(BookImage.south west)$)},
        y={($(BookImage.north west)-(BookImage.south west)$)}
      ]
        \ifstrequal{#2}{ch04-design.png}{%
          \node[diagramlabel,text width=.18\linewidth] at (.125,.125) {責務を分ける};
          \node[diagramlabel,text width=.18\linewidth] at (.375,.125) {境界を定める};
          \node[diagramlabel,text width=.18\linewidth] at (.625,.125) {変更を局所化};
          \node[diagramlabel,text width=.18\linewidth] at (.875,.125) {テストで検証};
        }{\ifstrequal{#2}{ch07-attention-generation.png}{%
          \node[diagramlabel,text width=.24\linewidth] at (.17,.15) {文脈を重み付きで参照};
          \node[diagramlabel,text width=.24\linewidth] at (.50,.15) {次トークンの確率を比較};
          \node[diagramlabel,text width=.24\linewidth] at (.83,.15) {選択して列へ追加};
        }{\ifstrequal{#2}{ch09-evidence-correlation.png}{%
          \node[diagramlabel,text width=.12\linewidth] at (.115,.76) {メトリクス\\傾向};
          \node[diagramlabel,text width=.12\linewidth] at (.115,.49) {ログ\\個別事象};
          \node[diagramlabel,text width=.12\linewidth] at (.115,.22) {トレース\\要求の経路};
          \node[diagramlabel,text width=.36\linewidth] at (.50,.075) {同じ時刻へそろえ、異常境界を絞る};
          \node[diagramlabel,font=\sffamily\bfseries\tiny,text width=.07\linewidth] at (.885,.78) {入口};
          \node[diagramlabel,font=\sffamily\bfseries\tiny,text width=.07\linewidth] at (.885,.62) {処理};
          \node[diagramlabel,font=\sffamily\bfseries\tiny,text width=.07\linewidth] at (.885,.44) {異常境界};
          \node[diagramlabel,font=\sffamily\bfseries\tiny,text width=.07\linewidth] at (.82,.27) {依存先A};
          \node[diagramlabel,font=\sffamily\bfseries\tiny,text width=.07\linewidth] at (.95,.27) {依存先B};
        }{%
          \node[diagramlabel,anchor=south,text width=.82\linewidth] at (.50,.025) {#3};
        }}}%
      \end{scope}
    \fi
  \end{tikzpicture}%
}

\newcommand{\chapterimage}[4][]{%
  \begin{figure}[htbp]
    \centering
    \IfFileExists{assets/#2}{%
      \BookAnnotatedImage{0.96\linewidth}{#2}{#1}%
    }{%
      \PackageError{cs-doc}{Required image `assets/#2' is missing}{Add the generated chapter image.}%
    }
    \caption{#3}
    \label{#4}
  \end{figure}
}

\newcommand{\sectionimage}[4][]{%
  \begin{figure}[htbp]
    \centering
    \IfFileExists{assets/#2}{%
      \BookAnnotatedImage{0.90\linewidth}{#2}{#1}%
    }{%
      \PackageError{cs-doc}{Required image `assets/#2' is missing}{Add the generated section image.}%
    }
    \caption{#3}
    \label{#4}
  \end{figure}
}

\newcommand{\takeaway}[1]{%
  \begin{takeawaybox}
  \textbf{要点}\quad #1
  \end{takeawaybox}
}

\newcommand{\source}[1]{%
  \par\smallskip
  {\raggedleft\footnotesize\color{PaleGray}出典：#1\par}%
  \medskip
}

\newcommand{\materialref}[3]{%
  \par\noindent\textbf{\href{#1}{#2}}\par
  \begingroup\small #3\par\endgroup
  \smallskip
}

\newcommand{\ChapterRef}[1]{第\ref{#1}章}
